\documentclass[10pt,a4paper]{amsart}
\usepackage{amsmath,amssymb,amsthm,graphicx}
\usepackage{verbatim}
\newenvironment{reference}{\comment}{\endcomment}
%\newenvironment{reference}{}{}
\newenvironment{slogan}{\comment}{\endcomment}
\newenvironment{history}{\comment}{\endcomment}
% For commutative diagrams we use Xy-pic
\usepackage[all]{xy}

% We use 2cell for 2-commutative diagrams.
\xyoption{2cell}
% Theorem environments.
%
\theoremstyle{plain}
\newtheorem{theorem}[subsection]{Theorem}
\newtheorem{proposition}[subsection]{Proposition}
\newtheorem{lemma}[subsection]{Lemma}

\theoremstyle{definition}
\newtheorem{definition}[subsection]{Definition}
\newtheorem{example}[subsection]{Example}
\newtheorem{exercise}[subsection]{Exercise}
\newtheorem{situation}[subsection]{Situation}

\theoremstyle{remark}
\newtheorem{remark}[subsection]{Remark}
\newtheorem{remarks}[subsection]{Remarks}

\numberwithin{equation}{subsection}

% Macros
%
\def\lim{\mathop{\mathrm{lim}}\nolimits}
\def\colim{\mathop{\mathrm{colim}}\nolimits}
\def\Spec{\mathop{\mathrm{Spec}}}
\def\Hom{\mathop{\mathrm{Hom}}\nolimits}
\def\Ext{\mathop{\mathrm{Ext}}\nolimits}
\def\SheafHom{\mathop{\mathcal{H}\!\mathit{om}}\nolimits}
\def\SheafExt{\mathop{\mathcal{E}\!\mathit{xt}}\nolimits}
\def\Sch{\mathit{Sch}}
\def\Mor{\mathop{\mathrm{Mor}}\nolimits}
\def\Ob{\mathop{\mathrm{Ob}}\nolimits}
\def\Sh{\mathop{\mathit{Sh}}\nolimits}
\def\NL{\mathop{N\!L}\nolimits}
\def\CH{\mathop{\mathrm{CH}}\nolimits}
\def\proetale{{pro\text{-}\acute{e}tale}}
\def\etale{{\acute{e}tale}}
\def\QCoh{\mathit{QCoh}}
\def\Ker{\mathop{\mathrm{Ker}}}
\def\Im{\mathop{\mathrm{Im}}}
\def\Coker{\mathop{\mathrm{Coker}}}
\def\Coim{\mathop{\mathrm{Coim}}}

% Boxtimes
%
\DeclareMathSymbol{\boxtimes}{\mathbin}{AMSa}{"02}

%
% Macros for moduli stacks/spaces
%
\def\QCohstack{\mathcal{QC}\!\mathit{oh}}
\def\Cohstack{\mathcal{C}\!\mathit{oh}}
\def\Spacesstack{\mathcal{S}\!\mathit{paces}}
\def\Quotfunctor{\mathrm{Quot}}
\def\Hilbfunctor{\mathrm{Hilb}}
\def\Curvesstack{\mathcal{C}\!\mathit{urves}}
\def\Polarizedstack{\mathcal{P}\!\mathit{olarized}}
\def\Complexesstack{\mathcal{C}\!\mathit{omplexes}}
% \Pic is the operator that assigns to X its picard group, usage \Pic(X)
% \Picardstack_{X/B} denotes the Picard stack of X over B
% \Picardfunctor_{X/B} denotes the Picard functor of X over B
\def\Pic{\mathop{\mathrm{Pic}}\nolimits}
\def\Picardstack{\mathcal{P}\!\mathit{ic}}
\def\Picardfunctor{\mathrm{Pic}}
\def\Deformationcategory{\mathcal{D}\!\mathit{ef}}
\usepackage{cleveref}
\crefname{originallemma}{original lemma}{original lemmas}
\Crefname{originallemma}{Original lemma}{Original lemmas}
\makeatletter
\newlabel{lemma-DM}{{lemma-DM}{source}{}{sourcepointer-1}{}}
\newlabel{lemma-DM@cref}{{[originallemma][0][]lemma-DM}{[1][0][]source}{}{}{}}
\makeatother
\crefname{originallemma}{original lemma}{original lemmas}
\Crefname{originallemma}{Original lemma}{Original lemmas}
\makeatletter
\newlabel{lemma-characterize-unramified}{{lemma-characterize-unramified}{source}{}{sourcepointer-2}{}}
\newlabel{lemma-characterize-unramified@cref}{{[originallemma][0][]lemma-characterize-unramified}{[1][0][]source}{}{}{}}
\makeatother
\crefname{originalreference}{original reference}{original references}
\Crefname{originalreference}{Original reference}{Original references}
\makeatletter
\newlabel{spaces-morphisms-lemma-etale-flat}{{spaces-morphisms-lemma-etale-flat}{source}{}{sourcepointer-3}{}}
\newlabel{spaces-morphisms-lemma-etale-flat@cref}{{[originalreference][0][]spaces-morphisms-lemma-etale-flat}{[1][0][]source}{}{}{}}
\makeatother
\crefname{originalreference}{original reference}{original references}
\Crefname{originalreference}{Original reference}{Original references}
\makeatletter
\newlabel{spaces-morphisms-lemma-etale-locally-finite-presentation}{{spaces-morphisms-lemma-etale-locally-finite-presentation}{source}{}{sourcepointer-4}{}}
\newlabel{spaces-morphisms-lemma-etale-locally-finite-presentation@cref}{{[originalreference][0][]spaces-morphisms-lemma-etale-locally-finite-presentation}{[1][0][]source}{}{}{}}
\makeatother
\crefname{originalreference}{original reference}{original references}
\Crefname{originalreference}{Original reference}{Original references}
\makeatletter
\newlabel{spaces-morphisms-lemma-etale-unramified}{{spaces-morphisms-lemma-etale-unramified}{source}{}{sourcepointer-5}{}}
\newlabel{spaces-morphisms-lemma-etale-unramified@cref}{{[originalreference][0][]spaces-morphisms-lemma-etale-unramified}{[1][0][]source}{}{}{}}
\makeatother
\crefname{originalreference}{original reference}{original references}
\Crefname{originalreference}{Original reference}{Original references}
\makeatletter
\newlabel{spaces-morphisms-lemma-unramified-flat-lfp-etale}{{spaces-morphisms-lemma-unramified-flat-lfp-etale}{source}{}{sourcepointer-6}{}}
\newlabel{spaces-morphisms-lemma-unramified-flat-lfp-etale@cref}{{[originalreference][0][]spaces-morphisms-lemma-unramified-flat-lfp-etale}{[1][0][]source}{}{}{}}
\makeatother
\title{Etaleness of a morphism of algebraic stacks tested through an etale presentation}
\author{The Stacks Project Authors}
\date{stacks-project:stacks-morphisms.tex:local-snapshot-20261012; GNU FDL 1.2 or later}
\begin{document}
\maketitle
\noindent\textit{Editorial note.} All delivered mathematical text is the author's own, copied from the pinned Stacks Project chapter source sources/stacks-algebraic-stacks-morphisms/pinned-original-source.tex; no statement or proof step was rewritten, reordered or removed.\par\medskip
\noindent\textit{Editorial note.} Preamble provenance: the Stacks Project keeps one single shared preamble for every chapter; this card takes it directly from the staged sources/stacks-algebraic-stacks-morphisms/preamble.tex as three verbatim spans (lines 8-12, 14-18 and 154-217), so no derived file was written and nothing was added or redefined. The wrapper class and the base packages are supplied by the assembler.\par\medskip
\noindent\textit{Editorial note.} Bibliography provenance: neither the counted statement nor its proof contains a citation (this lane rescans the whole span before writing the card), so the card carries no bibliography and no entry was reconstructed.\par\medskip
\noindent\textit{Editorial note.} Selection provenance: this card was opened under the file-based control contract (CONTROL.md), with the mechanical pre-card gate run on its statement and proof interval before the card was written.\par\medskip
\noindent\textit{Allowed standard prerequisites (editorial).}\begin{itemize}
\item The characterization of unramified morphisms used for one of the equivalences. Invoked at the start of the counted proof (source lines 8551-8553). Gives the equivalence of two of the listed conditions.
\item The characterization of DM morphisms through surjective etale covers. Invoked inside the counted proof (source lines 8555-8559). Produces the algebraic spaces presenting the morphism.
\item The statement that an unramified flat morphism of finite presentation of algebraic spaces is etale, together with the trivial implications from etaleness. Invoked at the end of the counted proof (source lines 8560-8566) as cross-chapter pointers. Finishes the equivalence for the presenting morphism.
\end{itemize}\medskip
\ifdefined\hypertarget\hypertarget{sourcepointer-1}{}\fi
\noindent\textit{Source reference.} Original lemma lemma-DM: sources/stacks-algebraic-stacks-morphisms/pinned-original-source.tex, source lines 4746--4746. Navigation pointer only: the counted text refers to the Stacks item lemma-DM, which is an existing labeled result of this same chapter and is not part of this card. The pointer carries the target's own label as its number and its definition line as its span; its content is not reproduced here.\par\medskip
\ifdefined\hypertarget\hypertarget{sourcepointer-2}{}\fi
\noindent\textit{Source reference.} Original lemma lemma-characterize-unramified: sources/stacks-algebraic-stacks-morphisms/pinned-original-source.tex, source lines 8450--8450. Navigation pointer only: the counted text refers to the Stacks item lemma-characterize-unramified, which is an existing labeled result of this same chapter and is not part of this card. The pointer carries the target's own label as its number and its definition line as its span; its content is not reproduced here.\par\medskip
\ifdefined\hypertarget\hypertarget{sourcepointer-3}{}\fi
\noindent\textit{Source reference.} Original reference spaces-morphisms-lemma-etale-flat: sources/stacks-algebraic-stacks-morphisms/pinned-original-source.tex, source lines 8566--8566. Navigation pointer only: the counted text refers to the Stacks item spaces-morphisms-lemma-etale-flat, which belongs to another chapter of the project and is not part of this card. The pointer carries the target's own label as its number because this chapter's staged snapshot does not contain the other chapter's tag table; its content is not reproduced here.\par\medskip
\ifdefined\hypertarget\hypertarget{sourcepointer-4}{}\fi
\noindent\textit{Source reference.} Original reference spaces-morphisms-lemma-etale-locally-finite-presentation: sources/stacks-algebraic-stacks-morphisms/pinned-original-source.tex, source lines 3672--3672. Navigation pointer only: the counted text refers to the Stacks item spaces-morphisms-lemma-etale-locally-finite-presentation, which belongs to another chapter of the project and is not part of this card. The pointer carries the target's own label as its number because this chapter's staged snapshot does not contain the other chapter's tag table; its content is not reproduced here.\par\medskip
\ifdefined\hypertarget\hypertarget{sourcepointer-5}{}\fi
\noindent\textit{Source reference.} Original reference spaces-morphisms-lemma-etale-unramified: sources/stacks-algebraic-stacks-morphisms/pinned-original-source.tex, source lines 8317--8317. Navigation pointer only: the counted text refers to the Stacks item spaces-morphisms-lemma-etale-unramified, which belongs to another chapter of the project and is not part of this card. The pointer carries the target's own label as its number because this chapter's staged snapshot does not contain the other chapter's tag table; its content is not reproduced here.\par\medskip
\ifdefined\hypertarget\hypertarget{sourcepointer-6}{}\fi
\noindent\textit{Source reference.} Original reference spaces-morphisms-lemma-unramified-flat-lfp-etale: sources/stacks-algebraic-stacks-morphisms/pinned-original-source.tex, source lines 1131--1131. Navigation pointer only: the counted text refers to the Stacks item spaces-morphisms-lemma-unramified-flat-lfp-etale, which belongs to another chapter of the project and is not part of this card. The pointer carries the target's own label as its number because this chapter's staged snapshot does not contain the other chapter's tag table; its content is not reproduced here.\par\medskip
\begin{lemma}
\label{lemma-characterize-etale}
Let $f : \mathcal{X} \to \mathcal{Y}$ be a morphism of algebraic stacks.
The following are equivalent
\begin{enumerate}
\item $f$ is \'etale, and
\item $f$ is locally of finite presentation, flat, and unramified,
\item $f$ is locally of finite presentation, flat, and its diagonal
is \'etale.
\end{enumerate}
\end{lemma}

\begin{proof}
The equivalence of (2) and (3) follows immediately from
Lemma \ref{lemma-characterize-unramified}. Thus in each case
the morphism $f$ is DM. Then we can choose Then we can choose
algebraic spaces $U$, $V$, a smooth surjective morphism
$V \to \mathcal{Y}$ and a surjective \'etale morphism
$U \to \mathcal{X} \times_\mathcal{Y} V$ (Lemma \ref{lemma-DM}).
To finish the proof we have to show that
$U \to V$ is \'etale if and only if it is locally of finite
presentation, flat, and unramified.
This follows from Morphisms of Spaces, Lemma
\ref{spaces-morphisms-lemma-unramified-flat-lfp-etale}
(and the more trivial
Morphisms of Spaces, Lemmas
\ref{spaces-morphisms-lemma-etale-unramified},
\ref{spaces-morphisms-lemma-etale-locally-finite-presentation}, and
\ref{spaces-morphisms-lemma-etale-flat}).
\end{proof}

\end{document}
